% answers_short.tex —— 1993 年·上海卷 · 简短答案速查
% 由 tools/make_short_answers.py 自 answers.tex 的 \gksolhead/\gkanswer 提取，
% 并与 1993/1993_上海_答案要点.md 逐题比对后生成，勿手改（改源 MD 后重跑）。

\begin{gknotice}
\textbf{使用说明}：本节只给最终答案便于快速核对，完整推导见解析版。答案为固定来源 PDF 卷内「解答」与评分标准所载（出版者/版次未核实），非官方标准答案；第 24、26 题的题面记号问题见答案要点注。
\end{gknotice}

\gkpart{填空题}
\begin{itemize}[leftmargin=2.6em]
  \item[\textbf{1.}] $2$。
  \item[\textbf{2.}] $\dfrac{\pi}{\omega}$。
  \item[\textbf{3.}] $\dfrac{3\pi}{5}$。
  \item[\textbf{4.}] $72$。
  \item[\textbf{5.}] $y=\cos x$，$\dfrac{\pi}{2}\le x\le\pi$。
  \item[\textbf{6.}] $2$。
  \item[\textbf{7.}] $600\pi$。
  \item[\textbf{8.}] $12$。
  \item[\textbf{9.}] $\rho\sin\theta=2$。
  \item[\textbf{10.}] $\arccos\dfrac{\sqrt3}{3}$。
\end{itemize}

\gkpart{选择题}
\begin{center}\begin{tabular}{ccccccccccc}
\toprule
题号 & 11 & 12 & 13 & 14 & 15 & 16 & 17 & 18 & 19 & 20 \\
\midrule
答案 & C & B & D & C & D & B & A & C & B & C \\
\bottomrule
\end{tabular}\end{center}

\gkpart{解答题}
\begin{itemize}[leftmargin=2.6em]
  \item[\textbf{21.}] $\dfrac{4-3\sqrt3}{10}$。
  \item[\textbf{22.}] $|w|=\dfrac{3\sqrt5}{2}$。
  \item[\textbf{23.}] 直线 $l$ 的方程为 $y=x-1$。
  \item[\textbf{24.}] （1）见解析中的证明；（2）$\dfrac{\sqrt3}{4}a$；（3）$CR=\dfrac{a}{2}$。
  \item[\textbf{25.}] （1）$a_n=a+2(n-1)b$，即 $\{a_n\}$ 是以 $a$ 为首项、$2b$ 为公差的等差数列；（2）直线方程为 $x-2y+a-2=0$；（3）$r\in(0,1)\cup\left(1,\dfrac52-\sqrt2\right)\cup(4+\sqrt6,+\infty)$。
  \item[\textbf{26.}] （1）$\operatorname{tg}\varphi=-\dfrac{2\sqrt{a^2+1}\,y_0}{(a^2-1)y_0^2+1}$（$y_0$ 为点 $P$ 的纵坐标）；（2）$\varphi_0=\pi-\operatorname{arctg}\sqrt{\dfrac{a^2+1}{a^2-1}}$；（3）$\dfrac{2\pi}{3}\le\varphi_0\le\pi-\operatorname{arctg}\sqrt2$。
\end{itemize}
